% A neural surrogate for the sodium-boiling phase of an SFR ULOF transient.
%
%   latexmk -pdf paper/paper.tex        (or: pdflatex paper/paper.tex, twice)
%
% LAYOUT: IOP Journal of Physics: Conference Series, from __DEV/paper_template.docx.
% The template is a Word *layout guide*, not a submission binary, so this file encodes
% its specification the same way `iop-jpcs.typ` does -- every constant below is taken
% from the guide's Table 1 (margins) and Table 2 (section styles), and the two files
% must be changed together. Deliberately NOT built on `iopart.cls`: the class imposes
% its own geometry, which would silently override the guide it is supposed to follow,
% and it is not installed everywhere. If IOP later ask for `iopart.cls` specifically,
% the body below transfers unchanged -- only this preamble is discarded.
%
% Content is kept in step with `paper.md` by hand. Where they disagree,
% `docs/axial_nn.md` holds the measurement record and both are wrong until reconciled
% against it.
\documentclass[11pt,a4paper]{article}

% --- Table 1 of the guide: A4, top 4.0, bottom 2.7, left 2.5, right 2.5 cm, and no
% --- headers, footers or page numbers -- production adds those.
\usepackage[a4paper,top=4.0cm,bottom=2.7cm,left=2.5cm,right=2.5cm]{geometry}
\pagestyle{empty}

% ORDER MATTERS HERE, and both orderings below were learned from a failed build.
%
% `amsmath` must load BEFORE `newtxmath`, and `amssymb` must not load AFTER it: newtxmath
% already provides the AMS symbols, so amssymb redefines \Bbbk and the run dies with
% "Command \Bbbk already defined". amssymb is simply dropped -- nothing in this paper uses
% a symbol it alone provides (\lVert and \rVert are amsmath).
\usepackage[T1]{fontenc}
\usepackage[utf8]{inputenc}
\usepackage{amsmath}

% Times at 11 pt, via `mathptmx` -- URW Nimbus Roman, a metrically exact Times clone,
% which is what the guide asks for ("11 point Times or Times New Roman").
%
% ONE font path, deliberately, after `newtx` cost this lane three red runs. Each was a
% different missing piece of the same package family: `binhex.tex` in
% texlive-plain-generic, then \Bbbk clashing with amssymb, then the TS1 metric
% `ts1-qtmr` from tex-gyre. Every one of them was invisible on the development machine,
% because newtx is not installed there and the conditional silently took the mathptmx
% branch -- so a green local build proved nothing about the branch CI actually ran.
%
% A conditional that selects an untestable path is worse than no conditional. mathptmx
% ships in texlive-fonts-recommended, is what the local build already used, and needs
% neither tex-gyre nor plain-generic. If newtx is ever wanted back, it needs
% texlive-fonts-extra, texlive-plain-generic and tex-gyre installed together, and amssymb
% must stay dropped.
\usepackage{mathptmx}
\usepackage{graphicx}
% Drawn by `uv run pinn-ulof figures`, never committed and never exported by hand --
% CI regenerates them from the physics before pdflatex runs. Draw them first for a
% local build; the directory is gitignored.
%
% Both paths are listed because TeX resolves them against the working directory and
% not against this file, so the build works from the repository root
% (`latexmk -pdf paper/paper.tex`, which is what CI runs) and from inside `paper/`
% (`latexmk -pdf paper.tex`) without the figures having to exist in two places.
\graphicspath{{figures/}{paper/figures/}}
% `url` rather than `hyperref`: the bibliography carries one web address and this typesets
% it with correct line breaking without altering layout or loading a package that conflicts
% with the font path above. It ships in texlive-latex-recommended.
\usepackage{url}
\usepackage{titlesec}
\usepackage[font=small,labelsep=period]{caption}
\usepackage{booktabs}

\setlength{\parindent}{5mm}   % first-line indent, per the guide
\setlength{\parskip}{0pt}

% --- Table 2: sections 11 pt bold, subsections 11 pt italic, one line space before and
% --- none after. Subsubsections run into the paragraph and end with a full stop.
\titleformat{\section}{\normalsize\bfseries}{\thesection.}{0.5em}{}
\titlespacing*{\section}{0pt}{\baselineskip}{0pt}
\titleformat{\subsection}{\normalsize\itshape}{\thesubsection}{0.5em}{}
\titlespacing*{\subsection}{0pt}{\baselineskip}{0pt}
\titleformat{\subsubsection}[runin]{\normalsize\itshape}{\thesubsubsection}{0.5em}{}[.]
\titlespacing*{\subsubsection}{0pt}{0.5\baselineskip}{0.4em}

% Captions: above tables, below figures, 10 pt. `caption`'s `font=small` gives 10 pt
% against an 11 pt body; the position is a matter of where \caption is placed.
\captionsetup[table]{position=top}
\captionsetup[figure]{position=bottom}
% "Captions should be below the figure and separated from it by a distance of 6 points."
% `singlelinecheck` is left ON deliberately: the guide centres short captions and
% justifies ones set to the width of the figure, which is exactly that default.
\setlength{\abovecaptionskip}{6pt}
\setlength{\belowcaptionskip}{0pt}
% "Table rules should be 0.5 points wide", and only three of them.
\setlength{\heavyrulewidth}{0.5pt}
\setlength{\lightrulewidth}{0.5pt}

% The guide: "a 5 mm gap between the reference number and the start of the reference
% text. Second and subsequent lines ... indented by 5 mm." `thebibliography` is used
% rather than a hand-numbered list so that \cite renumbers automatically -- inserting a
% section is exactly when a hand-numbered list goes wrong.
\makeatletter
\renewcommand\@openbib@code{%
  \setlength{\labelsep}{5mm}%
  \setlength{\itemsep}{0.4em}%
  \setlength{\parsep}{0pt}%
}
\makeatother

\begin{document}

\vspace*{28mm}    % "Leave 28 mm of space above the title and 10 mm after the title."

\begingroup
\raggedright
\fontsize{17}{20}\selectfont\bfseries
A neural surrogate for the sodium-boiling phase of an SFR unprotected loss-of-flow
transient: what it reproduces, and what it does not
\par
\endgroup

\vspace{10mm}

% `\raggedright` MUST come first: it resets \leftskip to zero, so setting the indent
% before it silently discards the 25 mm.
\begingroup
\raggedright \leftskip=25mm
M Lanzani\textsuperscript{1}
\par
\vspace{0.6em}
\fontsize{10}{12}\selectfont
\textsuperscript{1} Independent researcher\\
Corresponding author e-mail: marcolanzani03@gmail.com
\par
\endgroup

\vspace{1.2em}

% "indented 25 mm from the left margin" -- the left only.
\begingroup
\leftskip=25mm
\fontsize{10}{12}\selectfont
\noindent\textbf{Abstract.}
We solve the sodium-boiling phase of an unprotected loss-of-flow transient in a
sodium-cooled fast reactor with a physics-informed neural surrogate. One coolant channel is
resolved axially with four material fields and a sodium void fraction, coupled to six-group
point kinetics through Doppler and void-worth integrals, the latter changing sign near the
top of the core. Thermophysics, the boiling criterion and both feedback laws follow the
SAS4A/SASSYS-1 manual, with four registered deviations. Against a verified stiff Radau
reference the surrogate reaches a relative $L_2$ of $3.6\times10^{-4}$ on the film
temperature, places boiling onset 0.0131\,s from 10.9687\,s, and reproduces 100.2\% of the
peak voided length at a saturation margin of $+68.16$\,K against $+68.94$\,K. What the surrogate does not
yet deliver is the closed reactivity loop: the void-worth integral nearly cancels between
two large opposite contributions, and driving the kinetics from the learned fields recovers
only 8--16\% of it, while the non-cancelling Doppler integral is reproduced to within
1.7\%. We quantify that cancellation and show by an adjoint argument that it is not a
sampling deficiency.
\par
\endgroup

\vspace{10mm}

\section{Introduction}

In an unprotected loss-of-flow accident the primary pumps coast down while the reactor
protection system fails to scram. Coolant flow decays, the sodium outlet temperature rises,
and if saturation plus the required superheat is reached the coolant boils. In a
sodium-cooled fast reactor the resulting voiding carries a positive reactivity contribution
over most of the core height, and it competes with the prompt negative Doppler feedback.
Whether the excursion terminates benignly is decided by the sign and timing of that
competition, which makes the onset and axial extent of boiling the quantities of
engineering interest, not merely a field norm.

Whole-plant analysis of this sequence is the province of system codes, SAS4A/SASSYS-1 being
the reference of record \cite{sas4a}. A differentiable surrogate for the boiling phase is
attractive for uncertainty propagation, for parameter inference from limited
instrumentation, and for embedding a channel model inside an optimisation loop, all of
which require many evaluations and derivatives with respect to inputs.

This paper reports such a surrogate. Section~\ref{sec:model} states the governing equations
and closures, with the complete parameter set in table~\ref{tab:params}.
Section~\ref{sec:reference} describes the reference solution and quantifies its own
numerical uncertainty, which bounds every comparison that follows. Section~\ref{sec:method}
gives the surrogate. Section~\ref{sec:results} reports what it reproduces and what it does
not, and section~\ref{sec:discussion} discusses the consequences.

\section{Governing equations and closures}\label{sec:model}

\subsection{Channel model and state}

A single coolant channel is resolved along its axis, $\zeta = z/H \in [0,1]$, with radially
lumped materials. The state at each height comprises four temperatures --- fuel $T_f$,
cladding $T_{cl}$, structure $T_s$ and coolant $T_c$ --- and a void fraction $\alpha$,
together with six delayed-neutron precursor concentrations $c_i(t)$ that are functions of
time alone.

\subsection{Energy equations}

Writing $C_f$, $C_{cl}$ and $C_c$ for the heat capacities per unit channel length of fuel,
cladding and coolant, the field equations are
\begin{eqnarray}
C_f\,\partial_t T_f &=& (1-\gamma_c)\,q'(\zeta,t) - q_{fe}, \label{eq:fuel}\\
C_{cl}\,\partial_t T_{cl} &=& q_{fe} - q_{ec}, \label{eq:clad}\\
\rho_s c_s t_s\,\partial_t T_s &=& -\,h_{sc}(\alpha)\,(T_s - T_c), \label{eq:struct}\\
C_c\,\partial_t T_c &=& \bigl(1 - b(T_c)\bigr)\,q_w - w(t)\,c_c\,\partial_z T_c,
  \label{eq:cool}\\
\partial_t \alpha &=& S_v(T_c,\alpha,q_w) - u(t)\,\partial_z \alpha, \label{eq:void}
\end{eqnarray}
with $q_w = \gamma_c q' + q_{ec} + q_{sc}$ the total heat reaching the coolant and
$u = w/(\rho_c A_c)$ the coolant velocity. The linear power density is
$q'(\zeta,t) = P(t)\,P_0\,f(\zeta)/H$ with $f$ the axial power shape and $P$ the normalised
power of section~\ref{sec:kinetics}; a fraction $\gamma_c$ is deposited directly in the
coolant by neutron and gamma heating and the remainder in the fuel.

The inter-node fluxes carry both conduction and radiation across the gap,
\begin{equation}
q_{fe} = h_{gap} A_{fe} (T_f - T_{cl})
       + \sigma \varepsilon A_{fe} \bigl(T_f^4 - T_{cl}^4\bigr),
\end{equation}
while the wall-to-coolant paths are
$q_{ec} = h_{ec}(\alpha) A_{ec} (T_{cl} - T_c)$ and
$q_{sc} = h_{sc}(\alpha) A_{sc} (T_s - T_c)$. The film coefficients are interpolated
between their liquid and vapour values with the local void fraction, so that a voided node
insulates its wall.

Equations~(\ref{eq:fuel})--(\ref{eq:struct}) correspond to Eqs.~3.3-1 and 3.3-4 of the
manual, (\ref{eq:cool}) to Eq.~3.3-5 with its three source terms and the direct heating of
Eq.~3.3-6.

\subsection{Coolant transport and flow coastdown}

Pre-boiling momentum uses the manual's result (Eq.~3.9-1) that the mass flow is independent
of height, so $w = w(t)$ and the transport terms in (\ref{eq:cool}) and (\ref{eq:void}) are
purely advective. The coastdown is prescribed,
\begin{equation}
w(t) = w_0 \Bigl[ f_{\mathrm{nc}} + (1 - f_{\mathrm{nc}})\,
       e^{-t/\tau_{\mathrm{pump}}} \Bigr],
\end{equation}
with $f_{\mathrm{nc}}$ the natural-circulation floor.

\subsection{Sodium properties and the boiling criterion}

The eight relations of manual section~12.13 that this channel model evaluates are
implemented, Eq.~12.13-1 to 12.13-8: latent heat, saturation pressure and temperature,
liquid and vapour density, and liquid and vapour heat capacity. The remaining five ---
compressibility, thermal expansion, conductivity, viscosity and saturated liquid enthalpy
--- enter no term of the energy equations of section~\ref{sec:model} and are not carried.
Their stated validity is 590--2270\,K; the fits stop near 90\% of the critical point
(2503.3\,K) because $C_l$ is built from powers of $1/(T_c - T)$ and diverges there. The
solution is terminated when any temperature leaves that range, which defines the validity
horizon of section~\ref{sec:reference}.

Boiling begins where the coolant exceeds saturation by the required superheat,
\begin{equation}\label{eq:criterion}
T_c > T_{\mathrm{sat}}(p) + \Delta T_{\mathrm{sup}},
\end{equation}
following manual section~12.4, with the cladding-and-structure to vapour heat path of
section~12.5.1.

\subsection{Void closure}\label{sec:closure}

Rather than tracking vapour dynamics, the void fraction is slaved algebraically to the
coolant temperature,
\begin{equation}\label{eq:alpha}
\alpha = 1 - \bigl(1 - b(T_c)\bigr)^{3},
\end{equation}
where $b$ is a smooth switch that rises from zero to one over a width
$\Delta T_{\mathrm{sm}}$ about the criterion (\ref{eq:criterion}). The closure is justified
by timescale separation: vapour fills a node in 0.71\,ms against an advective 0.113\,s. It
was selected from six candidate forms on maximum voided-length error, at 1.2\% against the
reference's own 2.6\% mesh error at that resolution. The cubic exponent gives a finite
slope at $b = 0$, where a square-root closure has an unbounded one.

\subsection{Point kinetics and reactivity feedback}\label{sec:kinetics}

Power follows point kinetics with six delayed-neutron groups (Eq.~4.3-1), closed by the
prompt-jump approximation (Eq.~4.2-4). In normalised form, with $c_i(0) = P(0) = 1$,
\begin{equation}\label{eq:kinetics}
\frac{\mathrm{d}c_i}{\mathrm{d}t} = \lambda_i \bigl(P - c_i\bigr),
\qquad
P(t) = \frac{\sum_i \beta_i c_i}{\beta_{\mathrm{eff}} - \rho(t)} ,
\end{equation}
so the nominal state is an exact fixed point. The net reactivity is the sum of two axial
integrals, $\rho = \rho_D + \rho_{\mathrm{void}}$. The Doppler term is logarithmic in fuel
temperature and interpolated between flooded and voided coefficients (Eqs.~4.5-2, 4.5-3),
\begin{equation}\label{eq:doppler}
\rho_D(t) = \int_0^1 w_D(\zeta)\,\alpha_D(\alpha)\,
            \ln\frac{T_f(\zeta,t)}{T_f(\zeta,0)}\,\mathrm{d}\zeta ,
\end{equation}
and the coolant density and void are carried by a single worth distribution (Eq.~4.5-25),
\begin{equation}\label{eq:voidrho}
\rho_{\mathrm{void}}(t) = \int_0^1 w(\zeta)\,\alpha(\zeta,t)\,\mathrm{d}\zeta .
\end{equation}
The worth $w$ is positive over most of the core and negative near the top, changing sign at
$\zeta_{\mathrm{sign}}$ over a width $\delta_{\mathrm{sign}}$. That sign change is the
physical origin of the difficulty reported in section~\ref{sec:openloop}.

\subsection{Modelling assumptions}

Four departures from SAS4A are load-bearing and are stated here rather than left implicit.

\subsubsection{Prompt jump}
The prompt neutron mode is eliminated algebraically in (\ref{eq:kinetics}). This removes the
$\Lambda \approx 5\times10^{-7}$\,s timescale, which is the stiffest in the system. The
manual solves the full kinetics equations; this is a deviation, not a simplification of
them.

\subsubsection{Algebraic void closure}
Equation~(\ref{eq:alpha}) replaces vapour dynamics, as discussed in
section~\ref{sec:closure}. It is what allows a boiling front to form in a differentiable
formulation.

\subsubsection{Eulerian mixture and prescribed flow}
The void is treated as an Eulerian mixture rather than by Lagrangian slug tracking, and the
flow is prescribed rather than obtained from a pump-head balance.

Condensation, three-group decay heat and axial fuel expansion are implemented but disabled,
so no result depends on them; vapour expansion accelerating the flow is implemented but not
usable. Five further feedback mechanisms are omitted, and there is no pin failure or
material relocation, which bounds the transient to its pre-failure phase.

The complete parameter set --- every constant these equations use, with the value it was
solved at --- is given in appendix~\ref{app:params}.

\section{Reference solution and its numerical uncertainty}\label{sec:reference}

\subsection{Discretisation and integration}

The reference is a method-of-lines discretisation of
(\ref{eq:fuel})--(\ref{eq:void}) on a uniform axial mesh, with the advective terms treated
by first-order upwinding \cite{leveque}, integrated in time by the implicit Radau IIA
scheme \cite{hairer} as provided by SciPy \cite{scipy}. Radau is used because the system is
stiff even after the prompt-jump elimination: the fuel and cladding nodes relax on
millisecond timescales against an advective transit of order $0.1$\,s. The integration
supplies its own analytic Jacobian sparsity pattern.

It reproduces the analytic steady state to $3.4\times10^{-11}$\,K\,s$^{-1}$ and conserves
energy at second order in the step. Solution is terminated when any temperature reaches the
top of the section~12.13 property fits, which for these parameters occurs at 16.7\,s; the
surrogate is trained only over that horizon.

\subsection{Numerical uncertainty}

Because every claim below is stated against this solver, its own resolution bounds what can
be claimed. Refining the reference against itself in the manner of a grid-convergence study
\cite{eca} gives table~\ref{tab:ruler}.

\begin{table}[tbp]
\caption{The reference solution's own error at the scoring mesh, from mesh refinement.
These bound what any comparison against it can claim.}
\label{tab:ruler}
\centering
\begin{tabular}{@{}ll@{}}
\toprule
Quantity & Reference's own error \\
\midrule
Temperatures, relative $L_2$ & 1.1--$1.6\times10^{-3}$ \\
Temperature fields, relative $L_2$ & $9.7\times10^{-5}$ \\
Boiling onset time           & 0.00064\,s \\
Onset height                 & $0.0002$ in $\zeta$ \\
Peak voided length           & 0.049\% \\
Pointwise void fraction      & $2.0\times10^{-3}$ \\
\bottomrule
\end{tabular}
\end{table}

The estimates are obtained by solving on 640, 1280, 2560 and 5120 axial nodes and
extrapolating. The observed order of convergence is 1.00 on onset time and 1.00 on peak
voided length, consistent with the first-order upwinding of the advective terms, so
Richardson extrapolation applies and the uncertainty of a given mesh is its distance from
the extrapolated limit \cite{eca}. The difference between a mesh and the finest one
available is not its error, because that mesh is itself in error; the extrapolated limit is
the datum.

The saturation margin is excluded from that procedure. Its observed order is 0.70 rather
than 1, which is expected: it derives from a maximum over a discrete field rather than from
a smooth functional, so it is not in the asymptotic range and Richardson extrapolation is
not justified for it.

Refinement is cheap. The four solves take 8.3, 17.6, 31.6 and 61.1\,s, against hours to
train the surrogate they measure, so the resolution of the comparison is limited by choice
rather than by cost, and the comparisons of section~\ref{sec:results} are made against
the 2560-node solution.

Two consequences follow. An acceptance criterion on the pointwise void fraction is not
supportable, because the reference does not know that quantity to better than 3\%; such a
criterion was withdrawn. And calibration practice requires a tolerance to sit at least four
times above the uncertainty of the instrument measuring it \cite{nist}, so the 1\%
temperature criterion is sound at a ratio of about 100 and the 0.5\,s onset criterion at
about 780, but a result below that uncertainty is measuring the instrument
\cite{jakeman}.

Onset time is located by root-finding on the tangency of the level set
(\ref{eq:criterion}) rather than by inspection of the output grid. The reference onset is
10.9687\,s; read off a 0.25\,s output grid it appears as 10.75\,s, so a quarter of a second
of any onset error so obtained is quantisation.

\section{Surrogate formulation}\label{sec:method}

\subsection{Residual minimisation}

A physics-informed neural network represents the solution directly as a network $u_\theta$
of the independent variables and fits it by driving the residual of the governing equations
to zero at collocation points, using no solution data \cite{raissi,karniadakis}. For a
system written as $\mathcal{R}[u] = 0$ the objective is
\begin{equation}\label{eq:pinnloss}
\mathcal{L}(\theta) = \frac{1}{N}\sum_{i=1}^{N}
\bigl\lVert \mathcal{R}[u_\theta](\zeta_i, t_i)\bigr\rVert^2 ,
\end{equation}
with the derivatives that $\mathcal{R}$ requires obtained by automatic differentiation of
the network, exact to machine precision rather than differenced on a mesh. Which equations
enter $\mathcal{R}$, and with what weight, is not a matter of convenience and is set out in
section~\ref{sec:residualsystem}. The reference solution of section~\ref{sec:reference} is
used only to score the result, never to train it.

Three difficulties are established for this class of problem and each dictates a choice
below \cite{krishnapriyan}. The objective is a composition of the network with a
differential operator and is severely ill-conditioned, so a first-order method converges
only once its step size and moment parameters are tuned together, which fixes the optimiser
(section~\ref{sec:training}). Plain multilayer perceptrons learn low-frequency structure
long before high-frequency structure, which is fatal for a moving front and which fixes the
input embedding (section~\ref{sec:embedding}). And initial and boundary conditions added as
penalty terms compete with the residual through weights that must be chosen, which is
avoided here by imposing them exactly (section~\ref{sec:ansatz}).

\subsection{The residual system}\label{sec:residualsystem}

\subsubsection{Four residual blocks, not five}

The state has five fields, and section~\ref{sec:model} writes five field equations, but the
residual comprises only the four energy equations
(\ref{eq:fuel})--(\ref{eq:cool}). Equation~(\ref{eq:void}) is \emph{not} residualised, and
this is a consequence of the closure rather than an economy.

The closure (\ref{eq:alpha}) slaves $\alpha$ to $T_c$. Its justification is a timescale
separation --- vapour fills a node in 0.71\,ms against an advective 0.113\,s
(section~\ref{sec:closure}) --- and what that separation asserts is precisely that the void
transport equation is in quasi-steady balance on the timescale the model resolves. Once
$\alpha$ is a function of $T_c$, it is no longer an independent unknown: the system has four
unknown fields, and imposing five equations on them is over-determined. Worse, the fifth
equation is one the closure has deliberately approximated, so requiring the network to
satisfy (\ref{eq:void}) exactly asks it to undo the approximation that defines the model.

The optimiser has only one way to comply, which is to distort the temperature field, and it
does. Retaining the fifth block degrades the film temperature error by more than an order
of magnitude and places boiling onset seconds rather than milliseconds late; a five-block
solve at the full budget is worse than a four-block solve at one-fifteenth of it. The distinction matters for any formulation that eliminates a field algebraically:
\emph{the equation the closure replaces must leave the residual with it}, and a solver
that silently keeps both can converge smoothly to a wrong answer without any diagnostic
suggesting a problem.

\subsubsection{Variable scaling}

The four equations carry different units and magnitudes, and an unscaled sum lets the
stiffest field dominate. Each block is therefore multiplied by a constant $s_m$ chosen to
bring the four equations to a common magnitude. This changes the objective and never the
equations: the scales are constants, so dividing them back out recovers the physical
residual exactly and the zero set is untouched.

\subsubsection{Averaging over time windows}

The objective is not a plain mean over collocation points. The points are partitioned into
$K = 32$ equal windows $I_k$ in $\hat{t}$, each window averaged, and the windows then
averaged:
\begin{equation}\label{eq:objective}
\mathcal{L}(\theta) = \frac{1}{K}\sum_{k=1}^{K}\;
  \frac{1}{|I_k|}\sum_{i \in I_k}\; \sum_{m=1}^{4}
  \bigl[s_m\,r_m(\zeta_i,\hat{t}_i)\bigr]^2 .
\end{equation}

The reason is that a plain mean over points makes the objective depend on where the points
happen to lie. Under a plain mean, a choice about \emph{where to evaluate} the residual
silently becomes a choice about \emph{where to weight} it, and any density tilted toward
early time would weight it in exactly the wrong direction --- boiling onset occurs at
10.98\,s of a 16.5\,s window, so resolution of the coastdown would be bought at the expense
of the event the model exists to predict. Equation~(\ref{eq:objective}) keeps the two
decisions independent: every interval of the transient contributes equally to the loss
whatever the sampling does.

The sampling is uniform (section~\ref{sec:training}), so this is not a correction for a
deliberate density. It removes a Monte-Carlo one: a uniform draw fluctuates in how many
points land in each window, and (\ref{eq:objective}) makes the objective insensitive to
that fluctuation rather than letting it set the weights.

\subsection{Hard-constraint ansatz}\label{sec:ansatz}

The state is not the network output. It is written multiplicatively,
\begin{equation}\label{eq:ansatz}
\theta(\zeta,\hat{t}) = \theta_0(\zeta)\exp\bigl(\hat{t}\,N(\zeta,\hat{t})\bigr),
\end{equation}
with $\theta_0$ the analytic steady profile, $\hat{t} = t/t_{\mathrm{end}}$ and $N$ the
network. Three constraints then hold for any parameters: at $\hat{t} = 0$ the exponent
vanishes and the initial condition is the steady state to machine precision; the
exponential is positive, so no temperature falls below the inlet, which matters because the
Doppler term (\ref{eq:doppler}) is logarithmic in $T_f$; and a factor vanishing at
$\zeta = 0$ pins the single upstream boundary condition the advection admits. An additive
ansatz permits the optimiser to drive the fuel temperature negative while the loss falls,
at which point (\ref{eq:doppler}) is undefined.

\subsection{Embedding and network}\label{sec:embedding}

Figure~\ref{fig:network} shows the whole surrogate: two inputs, a frozen feature map, a
trunk, the ansatz that carries the constraints, and the residual blocks the objective is
built from. There is no encoder, no branch--trunk split and no solution data anywhere in it.

\begin{figure}[tbp]
\centering
\includegraphics[width=\linewidth]{network_diagram.png}
\caption{The surrogate, end to end. The Fourier matrix $B$ is drawn once and never trained,
so it is a change of input coordinates rather than a layer. Of the trainable parameters only
the trunk is fitting capacity; the read-out from the embedding grows with the number of
features and is listed separately in table~\ref{tab:netparams}. The ansatz makes the initial
and inlet conditions exact, so the objective contains residual terms only.}
\label{fig:network}
\end{figure}

The inputs pass through a random Fourier feature map \cite{tancik},
$x \mapsto [\sin(2\pi Bx), \cos(2\pi Bx)]$, with $B$ drawn once and held fixed: a change of
input coordinates rather than a fitted layer. The default uses 64 features, giving a
128-component embedding. The trunk is a 64-wide, 5-layer $\tanh$ multilayer perceptron with
five outputs, one per field. Only four of them are read: under the closure
(\ref{eq:alpha}) the void follows from $T_c$, and the fifth output is retained so that the
architecture does not depend on that modelling choice. Table~\ref{tab:netparams} separates
the parameters that fit from those that merely follow the embedding width.

\begin{table}[tbp]
\caption{Parameters of the default network. Only the trunk is fitting capacity; the
read-out width follows the embedding, and $B$ is frozen.}
\label{tab:netparams}
\centering
\begin{tabular}{@{}lr@{}}
\toprule
Component & Parameters \\
\midrule
Trunk (5 hidden layers, 64 wide, 5 outputs) & 17\,029 \\
Read-out from the embedding ($64 \times 128$) & 8\,192 \\
\midrule
Trainable total & 25\,221 \\
Frozen Fourier matrix $B$ (not trained) & 128 \\
\bottomrule
\end{tabular}
\end{table}

Embedding widths of 128 and 256 features were run at the same budget over three seeds. They
are indistinguishable from 64 on every quantity the reference can resolve, at roughly 1.25
and 1.9 times the cost per iteration, so the narrowest is used. A wider embedding does help
at a starved optimisation budget, where it improves onset accuracy and sharply reduces its
scatter across seeds; at the budget used here that advantage has gone.

\subsection{Collocation and training}\label{sec:training}

The collocation set is 10\,000 points drawn uniformly over
$(\zeta,\hat{t}) \in [0,1]^2$, once, and held fixed for the whole solve. There is no
sampling choice in the method: no density, no residual-adaptive refinement, no resampling.
Residual-adaptive refinement in particular is not used for the reason established in
section~\ref{sec:openloop}: once the void is closed algebraically the field residual is
small almost everywhere, so residual magnitude is a poor guide to where accuracy is needed.

The point count buys one quantity and only one. At 50\,000 iterations and three seeds,
boiling-onset error is 0.0314\,s at 5000 points, 0.0130\,s at 10\,000 and 0.0127\,s at
20\,000, so the improvement saturates at the value used here; over the same step the
half-range across seeds falls from 0.0197\,s to 0.0003\,s. The four temperature fields, the
peak voided length and the saturation margin converge to the same values at every point
count measured, which is the observation section~\ref{sec:discussion} returns to.

Training is 50\,000 L-BFGS iterations \cite{nocedal} with a strong-Wolfe line search and no
first-order stage, retaining 50 curvature pairs. All arithmetic is in double precision;
curvature pairs are formed from differences of gradients, and at single-precision residual
magnitudes those differences are dominated by rounding, so the quasi-Newton model would be
built from noise.

Each of those choices was measured. Adam-family optimisers were explored and dismissed:
they could reach comparable accuracy, but only after a time-consuming fine tuning of their
hyper-parameters. The collocation set is held fixed because the
curvature pairs a quasi-Newton method accumulates are meaningful only while the objective
is unchanged; redrawing during the solve invalidates that history and degrades the
converged result. At
10\,000 iterations no usable front forms at any point count, the peak voided length falling
short by more than 3\,cm and the saturation margin by 8\,K; the fields reach their converged
magnitudes by 30\,000 and onset stops improving measurably after 50\,000. The number of retained curvature pairs is a smaller
effect, 1.2 to 1.4 times in converged error between 10 and 50, but it is stated explicitly
because it is a common library default and leaving it unset once caused a difference in
convergence to be attributed to the software rather than to the argument.

\section{Results}\label{sec:results}

All results are at three seeds, against the reference at its scoring mesh of 2560 axial
nodes. Table~\ref{tab:ruler} is the output of the repository's \texttt{verify} command and
tables~\ref{tab:results} and \ref{tab:allseeds} of its \texttt{table} command; no value in
either is transcribed by hand.
Figure~\ref{fig:bc} shows the forcing that drives the sequence: the pumps coast down to the
natural-circulation floor while the power is still near nominal, so the outlet temperature
rises until it meets the criterion (\ref{eq:criterion}).

\begin{figure}[tbp]
\centering
\includegraphics[width=\linewidth]{boundary_conditions.png}
\caption{Prescribed flow coastdown and the coolant inlet and outlet response. The dashed
line marks boiling onset.}
\label{fig:bc}
\end{figure}

\subsection{Temperature fields and the boiling front}

Table~\ref{tab:results} compares the two solutions on the quantities of engineering
interest. Each entry is a value with its uncertainty: for the reference, the numerical
uncertainty of table~\ref{tab:ruler}; for the surrogate, the half-range over three seeds.
The complete per-seed record is in appendix~\ref{app:results}.

\begin{table}[tbp]
\caption{Reference and surrogate on the safety-relevant quantities. Reference
uncertainties are from mesh refinement; surrogate uncertainties are the half-range over
three seeds. The reference peak coolant temperature carries no refinement estimate, so no
uncertainty is quoted for it or for the margin derived from it.}
\label{tab:results}
\centering
\begin{tabular}{@{}lrr@{}}
\toprule
Quantity & Reference & Surrogate \\
\midrule
Boiling onset time (s)          & $10.9687 \pm 0.0006$ & $10.9818 \pm 0.0003$ \\
Onset height $\zeta^{*}$        & $0.9998 \pm 0.0002$  & $0.9998 \pm 0.0000$  \\
Peak voided length (m)          & $0.3785 \pm 0.0002$  & $0.3792 \pm 0.0002$  \\
Peak coolant temperature (K)    & $1237.9$             & $1237.1 \pm 0.1$     \\
Saturation margin (K)           & $+68.94 \pm 0.04$    & $+68.16 \pm 0.05$    \\
Film temperature, relative $L_2$& $9.7\times10^{-5}$   & $(3.64 \pm 0.09)\times10^{-4}$ \\
\bottomrule
\end{tabular}
\end{table}

On the four quantities where both uncertainties are known --- onset time, onset height,
peak voided length and the film temperature --- the surrogate agrees with the reference to
within their sum on the onset height alone. Onset time is 0.0131\,s late against a combined
0.0009\,s, and the peak voided length 0.0007\,m long against a combined 0.0004\,m. Both sit
outside the interval the two uncertainties allow, and at this mesh the reference contributes
only 0.0006\,s and 0.0002\,m of them, so both discrepancies are the model's and not the
ruler's. The seed half-ranges are small enough --- 0.0003\,s and 0.0002\,m --- that neither
can be attributed to the initialisation.

The film temperature is the quantity this reference cannot resolve. Its relative $L_2$ error
is $3.64\times10^{-4}$ against the reference's own $9.7\times10^{-5}$, a ratio of 3.8.
Calibration practice requires a tolerance to sit at least four times above the uncertainty
of the instrument measuring it; below that the two are not separable, and a smaller number
would be measuring the reference rather than the surrogate. The 1\% acceptance criterion is
met with three orders of magnitude to spare, but the size of what remains is not established
by this reference. It is not a matter of training longer: the converged film error is
$3.5\times10^{-4}$ at every point count and iteration budget measured, so separating it from
the discretisation error requires a finer reference than the convergence sequence of
table~\ref{tab:ruler} reaches.

Figures~\ref{fig:tmap} to \ref{fig:front} show what those scalars summarise: the coolant
field over space and time with the saturation contour that is also the front,
figure~\ref{fig:tmap}; the void fraction that follows from it, figure~\ref{fig:void}; and
the front height and voided length that the safety argument is stated in,
figure~\ref{fig:front}.

The peak coolant temperature is 0.8\,K low, and the saturation margin correspondingly
0.78\,K conservative. The reference carries no refinement estimate for either, so that
difference cannot be tested against an uncertainty and is reported as a difference rather
than as an agreement or a discrepancy.

Onset height is not a discriminating quantity in this model. The coolant heats
monotonically up the channel, so the hottest point, and therefore where boiling begins, is
always the outlet for surrogate and reference alike; a height criterion here measures the
mesh.

\begin{figure}[tbp]
\centering
\includegraphics[width=\linewidth]{temperature_map.png}
\caption{Coolant temperature over space and time. The contour is the criterion
(\ref{eq:criterion}), which under the closure (\ref{eq:alpha}) is the boiling front itself
rather than a separately computed curve; the marker is onset, located by tangency.}
\label{fig:tmap}
\end{figure}

\begin{figure}[tbp]
\centering
\includegraphics[width=\linewidth]{vapor_fraction.png}
\caption{Void fraction over space and time, with axial profiles at five instants and at
onset. The front forms at the outlet and propagates downward as the coolant heats.}
\label{fig:void}
\end{figure}

\begin{figure}[tbp]
\centering
\includegraphics[width=\linewidth]{front_height.png}
\caption{Front height and voided length. The saturation level set and the
$\alpha > 0.5$ contour do not coincide; the gap between them is the partially voided region
to which the worth integral (\ref{eq:voidrho}) is most sensitive.}
\label{fig:front}
\end{figure}

\subsection{Sensitivity to the collocation set}\label{sec:colloc}

Three collocation counts were run to convergence, three seeds each, and scored against the
same reference. Table~\ref{tab:colloc} compares them at a matched budget of
50\,000 iterations; appendix~\ref{app:ladders} gives every measured rung of all three.

\begin{table}[tbp]
\caption{Collocation count at a matched budget of 50\,000 iterations, three seeds each.
Field errors are relative $L_2$ over the whole space--time domain; uncertainties are the
half-range across seeds. Reference values are $0.3785 \pm 0.0002$\,m and
$+68.94 \pm 0.04$\,K.}
\label{tab:colloc}
\centering
\begin{tabular}{@{}lrrr@{}}
\toprule
Quantity & 5000 & 10\,000 & 20\,000 \\
\midrule
Fuel $T_f$ ($\times10^{-4}$)     & $16.55 \pm 0.56$ & $17.42 \pm 0.23$ & $17.39 \pm 0.47$ \\
Cladding $T_{cl}$ ($\times10^{-4}$) & $23.70 \pm 0.93$ & $24.73 \pm 0.24$ & $24.74 \pm 0.53$ \\
Film $T_s$ ($\times10^{-4}$)     & $3.49 \pm 0.16$  & $3.64 \pm 0.09$  & $3.70 \pm 0.08$ \\
Coolant $T_c$ ($\times10^{-4}$)  & $4.05 \pm 0.13$  & $3.95 \pm 0.03$  & $3.97 \pm 0.02$ \\
Onset error (s)                  & $0.0314 \pm 0.0197$ & $0.0131 \pm 0.0003$ & $0.0127 \pm 0.0011$ \\
Peak voided length (m)           & $0.3792 \pm 0.0002$ & $0.3792 \pm 0.0002$ & $0.3795 \pm 0.0003$ \\
Saturation margin (K)            & $+68.16 \pm 0.03$ & $+68.16 \pm 0.05$ & $+68.14 \pm 0.07$ \\
\bottomrule
\end{tabular}
\end{table}

Boiling onset is the only quantity the collocation count improves. It falls by a factor of
2.4 from 5000 to 10\,000 points and then stops, 20\,000 points buying a further 0.0003\,s
for twice the cost. The reduction in scatter is larger than the reduction in the value: the
half-range across seeds falls from 0.0197\,s to 0.0003\,s, so at 10\,000 points a single
run carries nearly the information of three, where at 5000 the initialisation is a
first-order term in the answer.

Nothing else responds. The four field errors differ by at most 6.2\% across a four-fold
range of collocation points, and the peak voided length and saturation margin agree to
0.0003\,m and 0.02\,K. The same holds along the budget axis: table~\ref{tab:lad5000},
table~\ref{tab:lad10000} and table~\ref{tab:lad20000} give the measured rungs, and from
40\,000 iterations upward every rung of every
count lies between 0.3791 and 0.3795\,m in peak voided length and between $+68.11$ and
$+68.18$\,K in margin, a spread of 0.0004\,m over a four-fold range of collocation points
and a two-and-a-half-fold range of budget. Below 40\,000 the rungs are still moving: at
30\,000 the same quantity spans 0.3787 to 0.3798\,m. Since that limit is reached from below at 5000 points and from
above at 20\,000, and does not move with either the point count or the budget, the
0.0007\,m by which it misses the reference is a property of the formulation rather than of
the discretisation of the residual or of the optimisation.

That distinction matters for how the ladders should be read. Intermediate budgets pass
\emph{through} the reference value on the way to the limit --- the peak voided length runs
0.3763, 0.3787, 0.3792\,m at 5000 points and 0.3788, 0.3798, 0.3794\,m at 20\,000 --- so a
budget selected because it agrees with the reference has been selected at a zero crossing
and will not reproduce. The converged value is quoted throughout. What the budget axis does
change, and what no choice of it removes, is examined in section~\ref{sec:open}.

\subsection{The transient in full}\label{sec:full}

The figures in this section are the reference solution, computed independently of the
surrogate and reproduced by it to the accuracy of section~\ref{sec:results}. They are
included because the safety argument rests on the sequence as a whole and not only on the
scalars of table~\ref{tab:results}.

The order is causal. Power stays near nominal while the flow decays,
figure~\ref{fig:power}, so every material heats, figure~\ref{fig:thist}, and the wall
fluxes that carry that heat into the coolant, figure~\ref{fig:flux}, drive the outlet
towards saturation until the margin closes, figure~\ref{fig:margin}. The state at the end
of the validity window, figure~\ref{fig:finalprof}, carries the kink that is the front.

\begin{figure}[tbp]
\centering
\includegraphics[width=\linewidth]{power.png}
\caption{Normalised power and the two reactivity feedbacks driving it. Power stays close to
nominal through the coastdown --- the transient is unprotected, so nothing scrams the core
--- and turns over only once voiding begins.}
\label{fig:power}
\end{figure}

\begin{figure}[tbp]
\centering
\includegraphics[width=\linewidth]{temperature_history.png}
\caption{Channel-average and peak temperature of each material against time. The fuel and
cladding track one another through the gap conductance; the coolant peak is the quantity
that meets saturation first.}
\label{fig:thist}
\end{figure}

\begin{figure}[tbp]
\centering
\includegraphics[width=\linewidth]{heat_flux.png}
\caption{The three wall-to-coolant heat paths against time. Their sum is $q_w$ in
(\ref{eq:cool}); the drop after onset is the voided node insulating its own wall through
the void-dependent film coefficients.}
\label{fig:flux}
\end{figure}

\begin{figure}[tbp]
\centering
\includegraphics[width=\linewidth]{saturation_margin.png}
\caption{Margin to saturation, $T_{\mathrm{sat}}+\Delta T_{\mathrm{sup}} - T_c$, against
time at several heights. Onset is where the uppermost curve reaches zero; the margin
elsewhere is what table~\ref{tab:results} reports as the saturation margin.}
\label{fig:margin}
\end{figure}

\begin{figure}[tbp]
\centering
\includegraphics[width=\linewidth]{final_temperature_profile.png}
\caption{Axial profile of all four material temperatures at the end of the validity window,
16.5\,s. The kink in the coolant profile at $\zeta \approx 0.62$ is the boiling front.}
\label{fig:finalprof}
\end{figure}

\subsection{The closed void-reactivity loop}\label{sec:openloop}

With the thermal fields prescribed the surrogate is accurate. Driving the kinetics
(\ref{eq:kinetics}) from the learned fields is a different matter, and it fails in a
specific way. Over the same fields and the same network the Doppler integral
(\ref{eq:doppler}) is reproduced to a factor of 1.017, while the void integral
(\ref{eq:voidrho}) recovers only 8--16\% of the reference value. The two differ in one
respect: the Doppler weight has one sign, and the void worth changes sign near the top of
the core (figure~\ref{fig:worth}).

\begin{figure}[tbp]
\centering
\includegraphics[width=0.78\linewidth]{void_worth_split.png}
\caption{The void worth $w(\zeta)$ of (\ref{eq:voidrho}), shaded by sign. Positive over most
of the core and negative near the top, so the functional is a difference of two large
contributions rather than a sum of small ones.}
\label{fig:worth}
\end{figure}

\begin{table}[tbp]
\caption{The void functional split at the sign change of the worth, at the instant it
peaks.}
\label{tab:cancel}
\centering
\begin{tabular}{@{}ll@{}}
\toprule
Contribution & Value \\
\midrule
Positive-worth region, $J^{+}$ & $+4.656\times10^{-4}$ \\
Negative-worth region, $J^{-}$ & $-1.695\times10^{-4}$ \\
Sum, $J$                       & $+2.962\times10^{-4}$ \\
Cancellation ratio, $|J|/(|J^{+}|+|J^{-}|)$ & 0.466 \\
\bottomrule
\end{tabular}
\end{table}

Table~\ref{tab:cancel} splits the functional at the sign change. A relative error
$\epsilon$ on each half becomes $2.1\epsilon$ on the sum, so the functional is an
ill-conditioned target by construction; reporting it as one number hides which half is
wrong, and the halves are reported separately here.

This is not a sampling deficiency. Dual-weighted-residual theory \cite{becker} gives the
error in a functional as $\langle \mathcal{R}[u_\theta], z^{*}\rangle$ to leading order,
with $z^{*}$ the solution of the adjoint problem sourced by the functional's derivative.
For the advective coolant operator of (\ref{eq:cool}) the adjoint runs backwards in $\zeta$
and can be evaluated in closed form, and the result is a step: the void slope in
(\ref{eq:alpha}) underflows to exactly zero wherever the coolant is subcooled, so $z^{*}$ is
constant over the lower 72\% of the channel and zero above it. Every point below the front
carries equal sensitivity and every point above it carries none. Residual-magnitude
sampling concentrates points at the front, precisely where the functional is insensitive,
so it cannot help; a uniform sampler is already near optimal.


\begin{figure}[tbp]
\centering
\includegraphics[width=\linewidth]{error_vs_residual.png}
\caption{Left: the surrogate's coolant-temperature error against the reference.
Right: the norm of the PDE residual (\ref{eq:objective}) that training minimises, on the
same axes and on a logarithmic scale. The white line is the reference boiling front. The
two fields are not the same shape: the residual is small and nearly uniform except on a
thin locus at the front, while the error occupies the whole voided region above it.}
\label{fig:errres}
\end{figure}


\begin{figure}[tbp]
\centering
\includegraphics[width=\linewidth]{surrogate_vs_reference.png}
\caption{Surrogate (dashed) over reference (solid) at four instants, for the coolant
temperature and the void fraction. The two are indistinguishable at plotting accuracy,
including across the front; figure~\ref{fig:errres} shows the difference that this view
cannot resolve.}
\label{fig:overlay}
\end{figure}

\begin{figure}[tbp]
\centering
\includegraphics[width=0.78\linewidth]{front_position.png}
\caption{Height of the boiling front against time, surrogate and reference. The front is
first resolved at onset near the channel outlet and descends as the voided region grows.}
\label{fig:frontpos}
\end{figure}

Figure~\ref{fig:overlay} shows the fields themselves and figure~\ref{fig:frontpos} the
front they define; neither separates the surrogate from the reference by eye.
Figure~\ref{fig:errres} makes the difference concrete for the thermal fields. The residual the
network minimises is small and nearly uniform across the domain, rising only on a thin
locus at the front, whereas the error against the reference occupies the entire voided
region behind it and peaks at 1.6\,K. Residual magnitude and error are therefore
uncorrelated in the region that matters, which is the sampling statement above expressed
pointwise: an adaptive scheme that places collocation where the residual is largest would
concentrate on the front, and the front is not where the error is.

This is a property of the formulation rather than of this particular fit. Closing the void
algebraically removes the one equation whose residual would be large across the voided
region, so what remains there is a set of energy equations that the network satisfies
accurately while the fields they are driven by have already drifted.

Evaluated open-loop on the surrogate's own field the functional is in fact accurate: the
positive half to 1.66--2.66\% and the negative half exactly, the latter because the
negative-worth region is fully voided in surrogate and reference alike and the integral is
then fixed by geometry. The deficit therefore arises in the closed loop, where
$\rho_{\mathrm{void}}$ feeds back into (\ref{eq:kinetics}) and the error compounds, and not
in the evaluation of the functional on a given field. Figure~\ref{fig:rho} shows the
reference reactivity split by mechanism.

\begin{figure}[tbp]
\centering
\includegraphics[width=0.78\linewidth]{reactivity.png}
\caption{Reactivity split by mechanism, in units of $\beta_{\mathrm{eff}}$. Reporting only
the net would hide which component carries the error.}
\label{fig:rho}
\end{figure}

\section{Discussion}\label{sec:discussion}

The surrogate reproduces the thermal-hydraulic phase of this transient to the resolution of
the reference solver, including the safety-relevant front quantities: when boiling starts,
how far it extends, and by what margin. For applications that consume those fields ---
uncertainty propagation over thermophysical parameters, inference of boundary conditions
from limited instrumentation, gradient-based studies of the coastdown --- the model is
usable, and its derivatives with respect to inputs are exact by construction.

For applications requiring the feedback loop closed it is not. The obstruction is not
accuracy in any ordinary sense, since the fields are right and the open-loop functional is
right, but the conditioning of a near-cancelling integral inside a feedback path. That is a
property of the physics rather than of the surrogate: a positive void worth over most of
the core against a negative worth near the top is what makes SFR void feedback interesting,
and it is what makes it a hard target. Any neural approach to sodium void feedback will meet
the same factor of 2.1.

The peak voided length is the sharpest instance of a limit that is not an accuracy limit.
It converges to between 0.3791 and 0.3795\,m against the reference's
$0.3785 \pm 0.0002$\,m, and it converges there from either side and from every
configuration tried: at 5000, 10\,000 and 20\,000 collocation points, every rung from
40\,000 iterations to 100\,000 lies in that interval, and every saturation margin in
$+68.11$ to $+68.18$\,K. Intermediate budgets pass through the
reference value on the way --- which is why a converged number, not an agreeing one, is what
should be quoted. Because the converged answer is independent of the collocation count and
of the optimisation budget, the remaining 0.0007\,m is a property of the formulation and not
of the fit; the algebraic void closure of section~\ref{sec:closure} is the natural
suspect, and testing that means changing the closure rather than training longer.

Two consequences for practice follow. A surrogate should be qualified against the
functionals it will be used for and not only against field norms: a relative $L_2$ of
$1.6\times10^{-3}$ on temperature coexists here with an 84--92\% miss on a reactivity
integral built from the same field. And the reference's own resolution must be quantified
before an acceptance criterion is set; two criteria in this work were set without that
check, one was withdrawn, and the temperature result has since reached the point where the
instrument rather than the model is the limit.

\section{Open questions}\label{sec:open}

\subsection{The objective and the accuracy metrics diverge}

The scalar objective (\ref{eq:objective}) falls monotonically for the whole solve. The
quantities the model is judged on do not, and after a point they move in opposite
directions. Table~\ref{tab:tradeoff} takes the two ends of the budget axis at 10\,000
collocation points.

\begin{table}[tbp]
\caption{Two ends of the budget axis at 10\,000 collocation points, three seeds. The
objective falls throughout; the fuel and cladding fields do not.}
\label{tab:tradeoff}
\centering
\begin{tabular}{@{}lrrr@{}}
\toprule
Quantity & 20\,000 iters & 100\,000 iters & change \\
\midrule
Fuel $T_f$ ($\times10^{-4}$)        & $14.96 \pm 0.52$ & $17.71 \pm 0.05$ & $+18\%$ \\
Cladding $T_{cl}$ ($\times10^{-4}$) & $21.30 \pm 0.63$ & $25.18 \pm 0.03$ & $+18\%$ \\
Film $T_s$ ($\times10^{-4}$)        & $7.61 \pm 0.58$  & $3.49 \pm 0.06$  & $-54\%$ \\
Coolant $T_c$ ($\times10^{-4}$)     & $6.81 \pm 0.14$  & $3.85 \pm 0.00$  & $-43\%$ \\
Onset error (s)                     & $0.0146 \pm 0.0132$ & $0.0115 \pm 0.0011$ & $-21\%$ \\
\bottomrule
\end{tabular}
\end{table}

Figure~\ref{fig:trajectories} puts every metric on one axis, at the shipped collocation
count, to show the shape of it. Each curve is an error against the reference divided by the
reference's own uncertainty in the same quantity, so the horizontal line at four is the
level below which a difference is no longer resolvable.

\begin{figure}[tbp]
\centering
\includegraphics[width=\linewidth]{metric_trajectories.png}
\caption{Every metric against the budget at 10\,000 collocation points, normalised by the
reference's own uncertainty in that quantity; bands are the half-range over three seeds.
The four temperature fields are solid, the three front quantities dashed. All seven improve
together to 20\,000 iterations. Past that the fuel and cladding errors rise to a plateau
near 20 while the film and coolant errors continue down to the resolution limit of the
reference, and the peak voided length passes through agreement rather than settling at it.
The scalar objective falls monotonically throughout.}
\label{fig:trajectories}
\end{figure}

Eighty thousand further iterations more than halve the film error and improve boiling onset
by a fifth, while making the fuel and cladding fields 18\% worse. The seed half-ranges at
100\,000 iterations are $5\times10^{-6}$ and $3\times10^{-6}$ in relative $L_2$, two orders
below the shift, so this is not scatter. The same non-monotonicity appears at every collocation count, and the size of it tracks how
early the minimum falls: the fuel field bottoms at 30\,000 iterations at 5000 points and
rises 10\% by 100\,000, at 20\,000 iterations at 10\,000 points and rises 18\%, and at
40\,000 iterations at 20\,000 points and rises 3\%.

The mechanism is visible in the structure of the objective. The four blocks are brought to a
common magnitude by fixed constants $s_m$ chosen once from the equations' characteristic
rates. That equalises the blocks at initialisation; it says nothing about their marginal
returns during the solve. Once the fuel and cladding residuals are small, further reduction
of the total is bought most cheaply in the film and coolant blocks, and the optimiser buys
it there --- including where doing so requires the solid fields to give ground. A single
scalar has one direction of improvement, and past a certain point that direction is not
aligned with any of the individual quantities.

Two consequences follow. Any budget quoted for this model is a choice of operating point on
that trade, not a point on a monotone convergence curve, and the choice must be made against
the quantities the surrogate will be used for. And the constants $s_m$ are the natural place
to intervene: a residual weighting that responds to how much each block is still improving
--- adapted from gradient norms, from the rate of change of each block, or from the
functionals the surrogate is to be qualified against --- would let the solve keep reducing
the film error without spending the fuel field to do it. Bounded adaptation is essential
here: unbounded per-block weights on this system run away and degrade every field. Whether
such a scheme moves the converged \emph{limits} or only the path to them is untested, and it
is the single most useful experiment this work leaves open.

\subsection{Each quantity has its own convergence rate}

The ladders show four distinct behaviours over one budget axis. The solid-material fields
reach their best value at 20\,000--40\,000 iterations and then degrade. The film and coolant
fields improve monotonically to the end of every ladder run. The peak voided length and the
saturation margin settle into an interval of 0.0004\,m and 0.07\,K by 40\,000 iterations and
stay inside it to 100\,000. Boiling onset falls at every rung from 30\,000 upward and has not clearly
saturated: at 10\,000 points it is still improving at the last rung, and at 20\,000 points it
is flat between 90\,000 and 100\,000 at 0.0108\,s. Below 30\,000 iterations none of the four
behaviours holds --- no quantity there is monotone in the budget.

No single stopping rule serves all four. Stopping when the objective flattens is the worst
of them, because the objective is still falling where three of the four have stopped
responding. A defensible protocol would monitor the functionals themselves, which requires
the reference during training and therefore a held-out validation transient rather than the
scored one.

\subsection{Limits that further training cannot reach}

The peak voided length converges to 0.3792\,m against the reference's
$0.3785 \pm 0.0002$\,m, from either side, at every collocation count and budget measured
(section~\ref{sec:colloc}). The residual gap is therefore in the formulation. The algebraic
void closure of section~\ref{sec:closure} is the natural suspect: it replaces the void
transport equation with a quasi-steady relation on the coolant temperature, exact in the
limit of instantaneous vapour redistribution, and a 0.2\% bias in the extent of the voided
region is the kind of error a quasi-steady closure would produce near the front. Testing
that means restoring the transport equation with a relaxation time and measuring whether the
limit moves --- not training longer.

Boiling onset never agrees with the reference at any configuration in this study. The
discrepancy narrows from 0.0314\,s to 0.0115\,s as the point count and budget increase, but
the seed scatter narrows faster, so the disagreement becomes \emph{more} significant rather
than less: at the shipped configuration it is 0.0131\,s against a combined uncertainty of
0.0009\,s. Whether it shares the void closure's origin, or is a separate consequence of
resolving a tangency condition on a locally flat field, is not resolved here.

The film temperature has reached the resolution of the reference. Its converged error is
$3.5\times10^{-4}$ at every configuration measured, against the reference's own
$9.7\times10^{-5}$ --- a ratio of 3.6 to 3.8, below the factor of four a tolerance requires
above the uncertainty of the instrument measuring it. No improvement to the surrogate can be
demonstrated on this field without first refining the reference, which the convergence
sequence of table~\ref{tab:ruler} does not currently reach.

\section{Conclusions}

A physics-informed neural surrogate for the sodium-boiling phase of an SFR unprotected
loss-of-flow transient, built on SAS4A/SASSYS-1 thermophysics and feedback laws with four
registered deviations, reproduces the reference temperature fields to $3.6\times10^{-4}$
relative $L_2$ on the film temperature and the peak voided length to 0.2\% at a saturation
margin of $+68.16$\,K against $+68.94$\,K. Boiling onset is placed 0.0131\,s late against
10.9687\,s, twenty times the reference's own uncertainty. Onset and voided length are both
resolved as genuine errors of the model; the voided length converges to the same value at
every collocation count and budget measured, which places it in the formulation rather than
in the fit.

The closed void-reactivity loop is not reproduced, recovering 8--16\% of the reference
integral while the non-cancelling Doppler integral over the same fields is reproduced to
within 1.7\%.
The cause is quantified as a cancellation ratio of 0.466, amplifying any error on either
half by a factor of 2.1, and is shown by an adjoint argument not to be a sampling
deficiency. Closing that loop is the outstanding problem for this model, and we expect it to
be the outstanding problem for neural surrogates of sodium void feedback generally.


\section*{Acknowledgments}

The SAS4A/SASSYS-1 documentation maintained by Argonne National Laboratory was the sole
source for the thermophysics and feedback laws used here.

\section*{Data availability statement}

The data that support the findings of this study are openly available at
\url{https://github.com/lanza2810/pinn-ulof}. No experimental or proprietary data were used.
The reference solution, every figure and every table in this paper are regenerated from the
governing equations and the parameters of appendix~\ref{app:params} by the code in that
repository: \texttt{reference} solves the reference transient, \texttt{verify} produces its
grid-convergence uncertainty, \texttt{train} fits the surrogate, \texttt{table} produces
tables~\ref{tab:results} and \ref{tab:allseeds}, \texttt{ladder} produces the budget ladders
of appendix~\ref{app:ladders}, and \texttt{figures} draws every figure. Training the shipped
configuration takes about three hours on eight CPU cores; the trained network behind the
model-dependent figures is included in the repository, so those figures can be redrawn
without repeating it.

\begin{thebibliography}{99}

\bibitem{sas4a} Fanning T H (ed) 2017 \emph{The SAS4A/SASSYS-1 Safety Analysis Code System}
ANL/NE-16/19 (Argonne, IL: Argonne National Laboratory). Section and equation numbers cited
here follow the maintained online edition ANL/NSE-SAS/5.8.1,
\url{https://sas-doc.nse.anl.gov/latest/}

\bibitem{leveque} LeVeque R J 2002 \emph{Finite Volume Methods for Hyperbolic Problems}
(Cambridge: Cambridge University Press)

\bibitem{hairer} Hairer E and Wanner G 1996 \emph{Solving Ordinary Differential Equations
II: Stiff and Differential-Algebraic Problems} 2nd edn (Berlin: Springer)

\bibitem{scipy} Virtanen P \emph{et al} 2020 SciPy 1.0: fundamental algorithms for
scientific computing in Python \emph{Nat. Methods} \textbf{17} 261

\bibitem{eca} E\c{c}a L and Hoekstra M 2014 A procedure for the estimation of the numerical
uncertainty of CFD calculations based on grid refinement studies
\emph{J. Comput. Phys.} \textbf{262} 104

\bibitem{nist} Taylor B N and Kuyatt C E 1994 \emph{Guidelines for Evaluating and Expressing
the Uncertainty of NIST Measurement Results} NIST Technical Note 1297

\bibitem{jakeman} Jakeman J D, Barba L A, Martins J R R A and O'Leary-Roseberry T 2026
Verification and validation for trustworthy scientific machine learning
\emph{Mach. Learn.: Sci. Technol.} \textbf{7} 025055

\bibitem{raissi} Raissi M, Perdikaris P and Karniadakis G E 2019 Physics-informed neural
networks: a deep learning framework for solving forward and inverse problems involving
nonlinear partial differential equations \emph{J. Comput. Phys.} \textbf{378} 686

\bibitem{karniadakis} Karniadakis G E, Kevrekidis I G, Lu L, Perdikaris P, Wang S and Yang L
2021 Physics-informed machine learning \emph{Nat. Rev. Phys.} \textbf{3} 422

\bibitem{krishnapriyan} Krishnapriyan A S, Gholami A, Zhe S, Kirby R M and Mahoney M W 2021
Characterizing possible failure modes in physics-informed neural networks
\emph{Adv. Neural Inf. Process. Syst.} \textbf{34} 26548

\bibitem{tancik} Tancik M, Srinivasan P P, Mildenhall B, Fridovich-Keil S, Raghavan N,
Singhal U, Ramamoorthi R, Barron J T and Ng R 2020 Fourier features let networks learn high
frequency functions in low dimensional domains \emph{Adv. Neural Inf. Process. Syst.}
\textbf{33} 7537

\bibitem{nocedal} Nocedal J and Wright S J 2006 \emph{Numerical Optimization} 2nd edn
(New York: Springer)

\bibitem{becker} Becker R and Rannacher R 2001 An optimal control approach to a posteriori
error estimation in finite element methods \emph{Acta Numerica} \textbf{10} 1

\end{thebibliography}

\appendix
\clearpage

\section{Complete parameter set}\label{app:params}

Table~\ref{tab:params} lists every constant used by
(\ref{eq:fuel})--(\ref{eq:voidrho}). Nothing outside this paper is required to reproduce
either the reference solution or the surrogate.

\begin{table}[tbp]
\caption{Complete parameter set. Sodium properties are evaluated from the manual's
section~12.13 correlations at the local temperature; the coolant values listed are the
reference state used for the constant-property terms.}
\label{tab:params}
\centering
\begin{tabular}{@{}llll@{}}
\toprule
Symbol & Description & Value & Unit \\
\midrule
\multicolumn{4}{@{}l}{\emph{Geometry}} \\
$H$          & active height                    & 1.0                  & m \\
$r_{fo}$     & fuel outer radius                & $3.60\times10^{-3}$  & m \\
$r_{ci}$     & cladding inner radius            & $3.75\times10^{-3}$  & m \\
$r_{co}$     & cladding outer radius            & $4.25\times10^{-3}$  & m \\
$A_c$        & coolant flow area                & $3.34\times10^{-5}$  & m$^2$ \\
$D_h$        & hydraulic diameter               & $5.0\times10^{-3}$   & m \\
$t_{s}$      & structure thickness              & $2.0\times10^{-3}$   & m \\
$\gamma_2$   & structure heat-transfer fraction & 0.5                  & --- \\
\midrule
\multicolumn{4}{@{}l}{\emph{Materials}} \\
$\rho_f, c_f$      & fuel                        & 9700, 330       & kg\,m$^{-3}$, J\,kg$^{-1}$K$^{-1}$ \\
$\rho_{cl}, c_{cl}$& cladding                    & 7800, 550       & kg\,m$^{-3}$, J\,kg$^{-1}$K$^{-1}$ \\
$\rho_s, c_s$      & structure                   & 7800, 550       & kg\,m$^{-3}$, J\,kg$^{-1}$K$^{-1}$ \\
$\rho_c, c_c$      & coolant at $T = 700$\,K     & 849.09, 1272.29 & kg\,m$^{-3}$, J\,kg$^{-1}$K$^{-1}$ \\
\midrule
\multicolumn{4}{@{}l}{\emph{Heat transfer}} \\
$h_{gap}$    & fuel-to-cladding gap conductance & $1.0\times10^{4}$    & W\,m$^{-2}$K$^{-1}$ \\
$\varepsilon$& gap emissivity                   & 0.7                  & --- \\
$h_{ec}$     & cladding-to-coolant, liquid      & $8.0\times10^{4}$    & W\,m$^{-2}$K$^{-1}$ \\
$h_{sc}$     & structure-to-coolant, liquid     & $8.0\times10^{4}$    & W\,m$^{-2}$K$^{-1}$ \\
$h_{v}$      & wall-to-vapour                   & $5.0\times10^{2}$    & W\,m$^{-2}$K$^{-1}$ \\
$\gamma_c$   & direct neutron and gamma heating & 0.02                 & --- \\
\midrule
\multicolumn{4}{@{}l}{\emph{Point kinetics, six groups}} \\
$\beta_{\mathrm{eff}}$ & total delayed fraction & $3.5\times10^{-3}$   & --- \\
$\Lambda$    & prompt generation time           & $5.0\times10^{-7}$   & s \\
$\lambda_i$  & decay constants                  & \multicolumn{2}{l}{0.0124, 0.0305, 0.111, 0.301, 1.14, 3.01 s$^{-1}$} \\
$\beta_i$    & group fractions ($\times10^{-4}$)& \multicolumn{2}{l}{1.157, 7.665, 6.858, 13.823, 4.027, 1.470} \\
\midrule
\multicolumn{4}{@{}l}{\emph{Reactivity feedback}} \\
$\alpha_D$   & Doppler, flooded                 & $-6.0\times10^{-3}$  & --- \\
$\alpha_D$   & Doppler, voided                  & $-4.0\times10^{-3}$  & --- \\
$\Delta k_{\mathrm{void}}$ & net void worth     & $+2.0\times10^{-3}$  & --- \\
$\zeta_{\mathrm{sign}}$ & worth sign-change height & 0.80             & --- \\
$\delta_{\mathrm{sign}}$ & sign-change width    & 0.05                 & --- \\
\midrule
\multicolumn{4}{@{}l}{\emph{Sodium and the boiling criterion}} \\
$p$          & system pressure                  & $1.01325\times10^{5}$& Pa \\
$\Delta T_{\mathrm{sup}}$ & required superheat  & 10.0                 & K \\
$\Delta T_{\mathrm{sm}}$  & closure smoothing width & 2.0              & K \\
\midrule
\multicolumn{4}{@{}l}{\emph{Transient}} \\
$P_0$        & nominal channel power            & $5.0\times10^{4}$    & W \\
$w_0$        & initial mass flow                & 0.25                 & kg\,s$^{-1}$ \\
$T_{in}$     & inlet temperature                & 628.0                & K \\
$\tau_{\mathrm{pump}}$ & coastdown time constant & 5.0                 & s \\
$f_{\mathrm{nc}}$ & natural-circulation floor   & 0.15                 & --- \\
$t_{\mathrm{end}}$ & nominal horizon            & 60.0                 & s \\
\bottomrule
\end{tabular}
\end{table}


\section{Per-seed measurements}\label{app:results}

Table~\ref{tab:allseeds} is the complete record behind table~\ref{tab:results}: all four
temperature fields and every front quantity, for each of the three seeds, together with
the seed mean and half-range.

\begin{table}[!ht]
\caption{Per-seed surrogate measurements. Relative $L_2$ errors are taken over the whole
space--time domain at the scoring mesh. The onset error is the absolute difference from
the reference onset of 10.9687\,s.}
\label{tab:allseeds}
\centering
\begin{tabular}{@{}lrrrr@{}}
\toprule
Quantity & Seed 0 & Seed 1 & Seed 2 & Mean $\pm$ half-range \\
\midrule
\multicolumn{5}{@{}l}{\emph{Relative $L_2$ error, $\times10^{-4}$}} \\
Fuel $T_f$        & 17.19 & 17.39 & 17.65 & $17.42 \pm 0.23$ \\
Cladding $T_{cl}$ & 24.49 & 24.65 & 24.96 & $24.73 \pm 0.24$ \\
Film $T_s$        & 3.554 & 3.734 & 3.560 & $3.644 \pm 0.090$ \\
Coolant $T_c$     & 3.988 & 3.971 & 3.919 & $3.954 \pm 0.034$ \\
\midrule
\multicolumn{5}{@{}l}{\emph{Boiling front}} \\
Onset time (s)             & 10.9820 & 10.9818 & 10.9815 & $10.9818 \pm 0.0003$ \\
Onset error (s)            & 0.0133  & 0.0131  & 0.0128  & --- \\
Onset height $\zeta^{*}$   & 0.9998  & 0.9998  & 0.9998  & $0.9998 \pm 0.0000$ \\
Peak voided length (m)     & 0.3791  & 0.3792  & 0.3794  & $0.3792 \pm 0.0002$ \\
Peak coolant temp.\ (K)    & 1237.1  & 1237.1  & 1237.2  & $1237.1 \pm 0.1$ \\
Saturation margin (K)      & $+68.11$& $+68.14$& $+68.21$& $+68.16 \pm 0.05$ \\
Maximum void fraction      & 1.0000  & 1.0000  & 1.0000  & $1.0000$ \\
\bottomrule
\end{tabular}
\end{table}


\section{Collocation-count ladders}\label{app:ladders}

Every finished rung of the three collocation counts, three seeds each, scored against the
same 2560-node reference. Field errors are relative $L_2$ over the whole space--time domain
in units of $10^{-4}$; uncertainties are the half-range across seeds. The reference values
are a peak voided length of $0.3785 \pm 0.0002$\,m and a saturation margin of
$+68.94 \pm 0.04$\,K.

\begin{table}[!ht]
\caption{5000 collocation points. This ladder is measured to 50\,000 iterations; the other
two reach 100\,000.}
\label{tab:lad5000}
\centering\footnotesize\setlength{\tabcolsep}{4pt}
\begin{tabular}{@{}rrrrrrrr@{}}
\toprule
Iters & $T_f$ & $T_{cl}$ & $T_s$ & $T_c$ & Onset err (s) & $L_{\mathrm{void}}$ (m) & Margin (K) \\
\midrule
10\,000 & $395.23 \pm 158.66$ & $563.81 \pm 213.26$ & $118.41 \pm 43.24$ & $120.80 \pm 44.39$ & $0.1401 \pm 0.0560$ & $0.3354 \pm 0.0146$ & $+53.22 \pm 5.67$ \\
20\,000 & $20.75 \pm 8.67$ & $34.52 \pm 14.81$ & $11.58 \pm 2.97$ & $10.69 \pm 3.41$ & $0.0408 \pm 0.0388$ & $0.3763 \pm 0.0031$ & $+67.51 \pm 0.60$ \\
30\,000 & $15.95 \pm 2.72$ & $22.39 \pm 3.73$ & $4.90 \pm 0.50$ & $5.03 \pm 0.69$ & $0.0438 \pm 0.0304$ & $0.3787 \pm 0.0006$ & $+68.00 \pm 0.18$ \\
40\,000 & $16.12 \pm 0.72$ & $22.98 \pm 0.99$ & $3.84 \pm 0.23$ & $4.28 \pm 0.16$ & $0.0392 \pm 0.0284$ & $0.3792 \pm 0.0003$ & $+68.11 \pm 0.04$ \\
50\,000 & $16.55 \pm 0.56$ & $23.70 \pm 0.93$ & $3.49 \pm 0.16$ & $4.05 \pm 0.13$ & $0.0314 \pm 0.0197$ & $0.3792 \pm 0.0002$ & $+68.16 \pm 0.03$ \\
60\,000 & $17.07 \pm 0.09$ & $24.38 \pm 0.25$ & $3.51 \pm 0.07$ & $3.94 \pm 0.04$ & $0.0293 \pm 0.0181$ & $0.3792 \pm 0.0002$ & $+68.15 \pm 0.02$ \\
70\,000 & $17.03 \pm 0.45$ & $24.46 \pm 0.51$ & $3.47 \pm 0.05$ & $3.91 \pm 0.09$ & $0.0275 \pm 0.0170$ & $0.3792 \pm 0.0001$ & $+68.16 \pm 0.02$ \\
80\,000 & $17.31 \pm 0.31$ & $24.72 \pm 0.42$ & $3.35 \pm 0.04$ & $3.88 \pm 0.07$ & $0.0240 \pm 0.0136$ & $0.3791 \pm 0.0001$ & $+68.17 \pm 0.02$ \\
90\,000 & $17.30 \pm 0.10$ & $24.76 \pm 0.07$ & $3.35 \pm 0.02$ & $3.88 \pm 0.03$ & $0.0219 \pm 0.0107$ & $0.3791 \pm 0.0000$ & $+68.16 \pm 0.02$ \\
100\,000 & $17.55 \pm 0.17$ & $25.04 \pm 0.08$ & $3.33 \pm 0.07$ & $3.86 \pm 0.02$ & $0.0202 \pm 0.0096$ & $0.3791 \pm 0.0001$ & $+68.15 \pm 0.01$ \\
\bottomrule
\end{tabular}
\end{table}

\begin{table}[!ht]
\caption{10\,000 collocation points, the shipped count. The 50\,000-iteration row is the
configuration of table~\ref{tab:results}.}
\label{tab:lad10000}
\centering\footnotesize\setlength{\tabcolsep}{4pt}
\begin{tabular}{@{}rrrrrrrr@{}}
\toprule
Iters & $T_f$ & $T_{cl}$ & $T_s$ & $T_c$ & Onset err (s) & $L_{\mathrm{void}}$ (m) & Margin (K) \\
\midrule
10\,000 & $225.16 \pm 99.15$ & $317.16 \pm 134.45$ & $66.96 \pm 26.19$ & $68.13 \pm 26.95$ & $0.0686 \pm 0.0670$ & $0.3615 \pm 0.0065$ & $+59.90 \pm 3.12$ \\
20\,000 & $14.96 \pm 0.52$ & $21.30 \pm 0.63$ & $7.61 \pm 0.58$ & $6.81 \pm 0.14$ & $0.0146 \pm 0.0132$ & $0.3778 \pm 0.0011$ & $+67.97 \pm 0.13$ \\
30\,000 & $17.40 \pm 0.94$ & $24.73 \pm 1.10$ & $4.42 \pm 0.12$ & $4.38 \pm 0.22$ & $0.0160 \pm 0.0008$ & $0.3795 \pm 0.0003$ & $+68.17 \pm 0.09$ \\
40\,000 & $18.24 \pm 0.45$ & $25.76 \pm 0.67$ & $3.91 \pm 0.22$ & $3.99 \pm 0.04$ & $0.0141 \pm 0.0033$ & $0.3794 \pm 0.0001$ & $+68.14 \pm 0.08$ \\
50\,000 & $17.42 \pm 0.23$ & $24.73 \pm 0.24$ & $3.64 \pm 0.09$ & $3.95 \pm 0.03$ & $0.0130 \pm 0.0003$ & $0.3792 \pm 0.0002$ & $+68.16 \pm 0.05$ \\
60\,000 & $17.51 \pm 0.29$ & $24.87 \pm 0.29$ & $3.49 \pm 0.07$ & $3.90 \pm 0.03$ & $0.0128 \pm 0.0013$ & $0.3793 \pm 0.0002$ & $+68.18 \pm 0.02$ \\
70\,000 & $17.57 \pm 0.06$ & $24.92 \pm 0.11$ & $3.35 \pm 0.11$ & $3.88 \pm 0.02$ & $0.0123 \pm 0.0019$ & $0.3792 \pm 0.0001$ & $+68.16 \pm 0.02$ \\
80\,000 & $17.59 \pm 0.16$ & $25.03 \pm 0.18$ & $3.46 \pm 0.07$ & $3.87 \pm 0.01$ & $0.0121 \pm 0.0008$ & $0.3793 \pm 0.0002$ & $+68.16 \pm 0.03$ \\
90\,000 & $17.69 \pm 0.07$ & $25.09 \pm 0.11$ & $3.48 \pm 0.05$ & $3.86 \pm 0.02$ & $0.0118 \pm 0.0011$ & $0.3792 \pm 0.0001$ & $+68.16 \pm 0.02$ \\
100\,000 & $17.71 \pm 0.05$ & $25.18 \pm 0.03$ & $3.49 \pm 0.06$ & $3.85 \pm 0.00$ & $0.0115 \pm 0.0011$ & $0.3792 \pm 0.0001$ & $+68.17 \pm 0.02$ \\
\bottomrule
\end{tabular}
\end{table}

\begin{table}[!ht]
\caption{20\,000 collocation points.}
\label{tab:lad20000}
\centering\footnotesize\setlength{\tabcolsep}{4pt}
\begin{tabular}{@{}rrrrrrrr@{}}
\toprule
Iters & $T_f$ & $T_{cl}$ & $T_s$ & $T_c$ & Onset err (s) & $L_{\mathrm{void}}$ (m) & Margin (K) \\
\midrule
10\,000 & $208.26 \pm 116.95$ & $288.33 \pm 159.87$ & $61.79 \pm 30.98$ & $62.71 \pm 31.67$ & $0.0834 \pm 0.0638$ & $0.3661 \pm 0.0091$ & $+60.59 \pm 4.07$ \\
20\,000 & $17.47 \pm 4.55$ & $25.45 \pm 6.55$ & $8.26 \pm 0.87$ & $6.82 \pm 0.46$ & $0.0090 \pm 0.0081$ & $0.3788 \pm 0.0017$ & $+68.07 \pm 0.26$ \\
30\,000 & $17.50 \pm 1.51$ & $24.95 \pm 2.06$ & $4.67 \pm 0.49$ & $4.30 \pm 0.27$ & $0.0157 \pm 0.0034$ & $0.3798 \pm 0.0004$ & $+68.13 \pm 0.07$ \\
40\,000 & $17.29 \pm 0.22$ & $24.49 \pm 0.21$ & $3.80 \pm 0.04$ & $4.06 \pm 0.03$ & $0.0130 \pm 0.0015$ & $0.3794 \pm 0.0001$ & $+68.15 \pm 0.06$ \\
50\,000 & $17.39 \pm 0.47$ & $24.74 \pm 0.53$ & $3.70 \pm 0.08$ & $3.97 \pm 0.02$ & $0.0127 \pm 0.0011$ & $0.3795 \pm 0.0003$ & $+68.14 \pm 0.07$ \\
60\,000 & $17.67 \pm 0.61$ & $25.14 \pm 0.76$ & $3.62 \pm 0.08$ & $3.88 \pm 0.07$ & $0.0119 \pm 0.0015$ & $0.3793 \pm 0.0001$ & $+68.16 \pm 0.02$ \\
70\,000 & $17.76 \pm 0.53$ & $25.18 \pm 0.74$ & $3.52 \pm 0.06$ & $3.87 \pm 0.07$ & $0.0117 \pm 0.0012$ & $0.3793 \pm 0.0002$ & $+68.17 \pm 0.00$ \\
80\,000 & $17.74 \pm 0.03$ & $25.17 \pm 0.04$ & $3.50 \pm 0.03$ & $3.87 \pm 0.01$ & $0.0113 \pm 0.0007$ & $0.3792 \pm 0.0001$ & $+68.16 \pm 0.00$ \\
90\,000 & $17.80 \pm 0.22$ & $25.28 \pm 0.27$ & $3.50 \pm 0.03$ & $3.84 \pm 0.01$ & $0.0108 \pm 0.0006$ & $0.3793 \pm 0.0001$ & $+68.18 \pm 0.01$ \\
100\,000 & $17.86 \pm 0.07$ & $25.35 \pm 0.05$ & $3.53 \pm 0.05$ & $3.84 \pm 0.00$ & $0.0108 \pm 0.0005$ & $0.3792 \pm 0.0001$ & $+68.17 \pm 0.01$ \\
\bottomrule
\end{tabular}
\end{table}

At 10\,000 iterations no collocation count produces a usable boiling front: the peak voided
length is short by 3 to 4\,cm and the saturation margin by 8 to 16\,K, with seed half-ranges
of several kelvin. That rung is a failure to converge rather than a less accurate solution,
and it is excluded from the comparisons in section~\ref{sec:colloc}.

\end{document}
